\documentclass[11pt]{article}
\usepackage[margin=1in]{geometry}
\usepackage[utf8]{inputenc}
\usepackage{textcomp}
\usepackage{amssymb}
\usepackage{longtable}
\usepackage{booktabs}
\usepackage{array}
\usepackage{hyperref}
\usepackage{titling}
\usepackage{parskip}
\setlength{\parindent}{0pt}
\hypersetup{colorlinks=true,linkcolor=blue!50!black,urlcolor=blue!50!black,citecolor=blue!50!black}

\title{Independent Replication Report --- Guan et al.\ 2022,\\
\emph{Olsenella intestinalis} sp.\ nov., isolated from cow feces\\
\large PMID 35689096 \textbullet{} DOI 10.1007/s00203-022-03017-2 \textbullet{} BVBRC-129}
\author{X-100 Replication Project (BVBRC set, rank 61)}
\date{2026-07-06}

\begin{document}
\maketitle

\begin{abstract}
\noindent\textbf{Verdict: PARTIAL.} Independent replication of Guan et al.\ (2022, \emph{Arch.\ Microbiol.} \textbf{204}:384) on the publicly deposited artifacts (RefSeq \texttt{GCF\_023276655.1}, 16S \texttt{NR\_181929.1}) confirms the paper's core taxonomic conclusion (BGYT1${}^{\mathrm{T}}$ is a distinct \emph{Olsenella}/\emph{Parafannyhessea} species) and its primary 16S neighbor claim ($\sim$98.4\% to \emph{O.\ umbonata}), and exposes several concrete discrepancies: the paper's abstract genome length (2{,}476{,}083 bp) contradicts its own body text (2{,}453{,}694 bp) and both NCBI records (2{,}453{,}694 bp), its OrthoANIu-reported ANI (76.8\%) is 3--7\,pp below what fastANI (80.83\%), skani (79.43\%) and hand-rolled reciprocal ANIb (83.36\%) report for the same genome pair, and the paper's chitinase / $\beta$-1,3-glucanase claims are not supported by the current PGAP annotation. Phenotypic/chemotaxonomy claims are out-of-reach without the live isolate KCTC 25379${}^{\mathrm{T}}$.
\end{abstract}

\section{Paper summary}

Guan and co-authors describe a single novel bacterium, strain \textbf{BGYT1${}^{\mathrm{T}}$} ($=$ KCTC 25379${}^{\mathrm{T}}$ $=$ GDMCC 1.3011${}^{\mathrm{T}}$), isolated from cow feces in Andong, Republic of Korea. They use polyphasic taxonomy: 16S rRNA phylogeny (MEGA v7.0 NJ/ML/ME on Clustal W alignment), whole-genome sequencing (Illumina Nexa + Oxford Nanopore MinION, assembled with SPAdes v3.13.1), TYGS/GBDP phylogenomics with FastME v2.1.6.1, plus wet-lab morphology, physiology, biochemistry (API 20A/ZYM/Rapid ID 32A), fatty-acid profiling (MIDI + Sherlock Anaerobic v6.1) and polar-lipid TLC. Based on the combined evidence, they propose \emph{Olsenella intestinalis} sp.\ nov.\ with BGYT1${}^{\mathrm{T}}$ as the type strain.

The BV-BRC pipeline analog is \emph{Codon Tree / Phylogenetic Tree} on the newly-sequenced genome plus \emph{Olsenella} reference type strains --- i.e.\ a genome-based phylogenetic placement of a novel-species announcement.

\section{Claims table}

\begin{longtable}{@{}p{0.5cm}p{6.6cm}p{2cm}p{2cm}p{3.2cm}@{}}
\toprule
\# & Claim (paper) & Type & Tested? & Verdict \\
\midrule\endhead
C1  & Genome length: 2{,}476{,}083 bp (abstract) / 2{,}453{,}694 bp (body) & Quantitative & Yes & \textbf{REPLICATED} body; \textbf{CONTRADICTS} abstract (paper-internal error) \\
C2  & G+C = 66.9 mol\% & Quantitative & Yes & \textbf{REPLICATED} (66.95\% own count) \\
C3  & 8 contigs & Quantitative & Yes & \textbf{PARTIAL}; NCBI now 2 contigs \\
C4  & N50 = 604{,}117 bp & Quantitative & Yes & \textbf{PARTIAL}; NCBI 2-contig N50 = 1{,}425{,}513 \\
C5  & 1835 genes / 1778 CDS / 50 tRNAs / 6 rRNAs / 1 tmRNA & Quantitative & Yes & \textbf{REPLICATED} within re-annotation drift (1810/1761/49/6/1 in PGAP v6.11) \\
C6  & 16S closest neighbor = \emph{O.\ umbonata} KCTC 15140${}^{\mathrm{T}}$ at 98.24\% & Quantitative & Yes & \textbf{REPLICATED} (98.38\% own pairwise) \\
C7  & ANI to closest = 76.8\% & Quantitative & Yes & \textbf{PARTIAL} (fastANI 80.83, skani 79.43, ANIb 83.36; qualitatively agrees, magnitude differs) \\
C8  & AAI to \emph{O.\ umbonata} = 67.3\% & Quantitative & No & \textbf{NOT\_TESTED} (deferred) \\
C9  & dDDH to \emph{O.\ umbonata} = 22.2\% & Quantitative & No & \textbf{NOT\_TESTED} (GGDC web-queued) \\
C10 & Novel species (below all species thresholds) & Qualitative & Yes & \textbf{REPLICATED} unambiguously \\
C11 & Placed in genus \emph{Olsenella}, sister to \emph{O.\ umbonata} & Qualitative & Yes & \textbf{REPLICATED} (own NJ tree over 13 type strains) \\
C12 & Chitinase / $\beta$-1,3-glucanase / protease genes detected & Presence & Yes & \textbf{PARTIAL} (proteases: 54 hits; chitinase: 0; glucanase: 0 in current PGAP) \\
C13 & Gram-negative, anaerobic, non-motile, rod-shaped & Phenotypic & No & \textbf{OUT\_OF\_REACH} (paper's own genus description says Actinobacteria are Gram-positive; needs isolate) \\
C14 & pH 6--8 (opt 7), 35--40\,\textdegree C (opt 35), 0.5--1.5\% NaCl (opt 0.5) & Phenotypic & No & \textbf{OUT\_OF\_REACH} \\
C15 & Dominant FAs: C\textsubscript{16:0} DMA 20.2, C\textsubscript{16:0} 20.2, C\textsubscript{18:0} 10.5, C\textsubscript{18:1} \emph{cis} 9 17.0\% & Chemotaxonomy & No & \textbf{OUT\_OF\_REACH} \\
C16 & Polar lipids: 1 phospholipid + 4 glycolipids + 3 lipids (all unidentified) & Chemotaxonomy & No & \textbf{OUT\_OF\_REACH} \\
\bottomrule
\end{longtable}

\textbf{Coverage:} 12 of 16 claims are testable from public data; 10 tested; 6 fully replicated; 4 partially; 1 (C1-abstract) directly contradicted by the paper's own body text; 4 out-of-reach (wet-lab only).

\section{Method}

\subsection{Data sources}
\begin{itemize}
\item \textbf{BGYT1 assembly:} NCBI Datasets \texttt{GCF\_023276655.1} (RefSeq, ASM2327665v1; released 2022-05-08 by KRIBB); paired \texttt{GCA\_023276655.1}.
\item \textbf{BGYT1 16S rRNA:} \texttt{NR\_181929.1} (RefSeq; original submission \texttt{OM533390.1}).
\item \textbf{Closest relative assembly:} \texttt{GCF\_900105025.1} \emph{Parafannyhessea umbonata} DSM 22620 (formerly \emph{Olsenella umbonata} KCTC 15140${}^{\mathrm{T}}$ = A2, reclassified by Zgheib et al.\ 2021).
\item \textbf{Closest relative 16S:} \texttt{AJ251324.3} (\emph{O.\ umbonata} A2, 1433 nt).
\item \textbf{Type-strain 16S panel (13 sequences)} for genus-scale phylogeny: NR\_181929 (BGYT1), AJ251324 (umbonata A2), NR\_199489 (absiana), NR\_180810 (porci), NR\_180580 (lakotia), NR\_179615 (congonensis), NR\_179506 (timonensis), NR\_173694 (massiliensis), NR\_173693 (phocaeensis), NR\_173692 (urininfantis), NR\_116939 + NR\_116938 (profusa), NR\_115110 (uli).
\item \textbf{Paper full text:} Springer paywalled; publicly-served article HTML page contains full body text and was used for extraction; PDF (\texttt{paper.pdf}) is a Chrome-headless render of that HTML.
\end{itemize}

\subsection{Tools and versions}
\begin{itemize}
\item \texttt{fastANI} (Homebrew), \texttt{skani} (Homebrew), BLAST+ 2.16 (\texttt{blastn}, \texttt{makeblastdb}); \texttt{clustalo}; Biopython 1.87 (\texttt{Bio.Align.PairwiseAligner}, \texttt{Bio.Phylo.TreeConstruction}); Python 3.14; \texttt{curl} 8.x for NCBI E-utilities and NCBI Datasets v2alpha REST; Chrome headless for HTML$\to$PDF.
\item MUMmer \texttt{nucmer}/\texttt{dnadiff} \emph{attempted} but skipped due to a Homebrew \texttt{TIGR::Foundation} perl-module packaging bug; ANIm was replaced by hand-rolled reciprocal ANIb.
\end{itemize}

\subsection{Commands (condensed)}
\begin{verbatim}
# Metadata
curl -s '.../efetch.fcgi?db=pubmed&id=35689096&retmode=xml' > pubmed.xml
curl -s '.../dataset_report' > assembly_report.json

# Genomes
curl -s -o BGYT1_genome.zip '.../GCF_023276655.1/download?include_annotation_type=GENOME_FASTA'
curl -s -o umbonata_genome.zip '.../GCF_900105025.1/download?include_annotation_type=GENOME_FASTA'

# Own contig / GC / N50 recount
python3 count_stats.py bgyt1.fna    # -> 2,453,694 bp / 66.95% GC / 2 contigs / N50 = 1,425,513

# 16S pairwise (Biopython local aligner, +1/-1/-2/-0.5)
python3 pairwise_16s.py             # -> 98.38%

# ANI (three independent methods)
fastANI -q bgyt1.fna -r umbonata.fna -o fastani.tsv                 # 80.83%
fastANI -q umbonata.fna -r bgyt1.fna -o fastani_rev.tsv             # 80.76%
skani dist bgyt1.fna umbonata.fna -s 70 --slow                       # 79.43%
python3 anib_reciprocal.py          # 1020 bp frags, blastn, >=30% id, >=70% cov -> 83.36%

# Phylogeny (13 type strains)
clustalo -i olsenella_type_strains_16S.fasta -o msa.fasta --force
python3 nj_tree.py                  # Biopython DistanceCalculator('identity') + NJ

# Annotation cross-check
awk -F'\t' '!/^#/ {print $3}' genomic.gff | sort | uniq -c
grep -Ei 'chitinase|glucanase|protease|peptidase' genomic.gff
\end{verbatim}

\section{Results vs.\ paper}

\subsection{Genome metrics}
\begin{tabular}{@{}lrrrrl@{}}
\toprule
Metric & Abstract & Body & NCBI & Own count & Verdict \\
\midrule
Length (bp) & 2{,}476{,}083 & 2{,}453{,}694 & 2{,}453{,}694 & 2{,}453{,}694 & Body replicated; abstract contradictory \\
GC\% & 66.9 & 66.9 & 67 & 66.95 & Replicated \\
Contigs & 8 & 8 & 2 & 2 & NCBI cleanup \\
N50 (bp) & 604{,}117 & 604{,}117 & 1{,}425{,}513 & 1{,}425{,}513 & Reflects contig-count difference \\
Genes & 1835 & 1835 & 1810 & --- & Re-annotation drift $<$2\% \\
CDS & 1778 & 1778 & 1761 & --- & Re-annotation drift $<$1\% \\
tRNA & 50 & 50 & 49 & --- & 1 fewer \\
rRNA & 6 & 6 & 6 & --- & Exact (2$\times$ 5S/16S/23S) \\
tmRNA & 1 & 1 & 1 & --- & Exact \\
\bottomrule
\end{tabular}

\subsection{16S neighbor identity}
Own Biopython pairwise (local, $+1/{-1}/{-2}/{-0.5}$) BGYT1 vs.\ \emph{O.\ umbonata} A2 = \textbf{98.38\%} (matches / aln length) or 98.45\% (matches / (matches+mismatches)). Paper (EzBioCloud OrthoANIu-family): 98.24\%. Difference $<$0.2 pp --- excellent agreement.

Own full-genus 16S similarity ranking (top 6 hits):

\begin{tabular}{@{}lr@{}}
\toprule
Neighbor & \% id (aln) \\
\midrule
\emph{Parafannyhessea umbonata} (was \emph{O.\ umbonata}) A2 (AJ251324.3) & \textbf{98.38} \\
\emph{Olsenella profusa} DSM 13989 (NR\_116939.1) & 97.68 \\
\emph{Olsenella phocaeensis} Marseille-P2936 (NR\_173693.1) & 97.54 \\
\emph{Olsenella timonensis} Marseille-P2300 (NR\_179506.1) & 97.19 \\
\emph{Olsenella lakotia} SW165 (NR\_180580.1) & 96.84 \\
\emph{Olsenella uli} DSM 7084 (NR\_115110.1) & 96.76 \\
\bottomrule
\end{tabular}

Rank order agrees with paper (\emph{umbonata} $>$ \emph{profusa} $>$ \emph{uli}). Our NJ tree over 13 type strains (\texttt{report/evidence/olsenella\_16S\_NJtree.newick}) places BGYT1 as sister to \emph{P.\ umbonata}, with \emph{O.\ profusa} as the next-closest branch --- matches paper's Fig.\ 1.

\subsection{ANI vs.\ closest relative}
\begin{tabular}{@{}lrrl@{}}
\toprule
Method & Result & Paper & Interpretation \\
\midrule
fastANI (BGYT1$\to$umb) & 80.83\% & 76.8 & above paper \\
fastANI (umb$\to$BGYT1) & 80.76\% & 76.8 & reciprocal-consistent \\
skani (learned, $-s\ 70$ \texttt{--slow}) & 79.43\% & 76.8 & above paper \\
Own reciprocal ANIb (1020 bp frags, $\geq$30\% id, $\geq$70\% cov) & 83.36\% & 76.8 & above paper \\
\midrule
\textbf{All methods:} $\ll$ 95\% species threshold & --- & --- & \textbf{species novelty REPLICATED} \\
\bottomrule
\end{tabular}

The magnitude gap ($\approx$3--7 pp above paper for every modern tool) is systematic and consistent with an OrthoANIu-vs-modern-tools bias at cross-genus divergences. This does not affect the taxonomic conclusion.

\subsection{Annotation feature cross-check}
Current PGAP v6.11 (2026-05-18) on \texttt{GCF\_023276655.1}: 1810 genes / 1761 CDS / 49 tRNAs / 6 rRNAs / 1 tmRNA / 9 pseudogenes / 4 riboswitches. Peptidase/protease grep: 54 hits (supports paper's protease claim). Chitinase grep: \textbf{0 hits}. Explicit $\beta$-1,3-glucanase grep: \textbf{0 hits}. GH families found: GH1 (2), GH3 (2), GH25 (2, lysozyme-like), plus 3 $\alpha$-amylase-family GH --- bacterial-wall / carbohydrate machinery, not fungal-wall-degrading machinery. The paper's chitinase/glucanase-based biocontrol motivation is not supported by the current annotation.

\section{Per-claim what-worked / what-didn't}

\begin{longtable}{@{}p{1cm}p{5.5cm}p{9cm}@{}}
\toprule
Claim & What worked & What did not \\
\midrule\endhead
C1  & Body-text length 2{,}453{,}694 bp reproduced exactly from own contig-sum, matches both NCBI records & Abstract's 2{,}476{,}083 bp is contradicted by every source, including the paper's own body \\
C2  & 66.95\% GC (own) matches paper's 66.9\% & --- \\
C3--C4 & --- & Current NCBI records show 2 contigs / N50 1.4 Mb, not paper's 8 / 604 kb; either NCBI cleaned/merged small contigs or the paper's SPAdes+MinION assembly was later polished \\
C5  & Feature counts within 1--2\% of paper (drift band for 4-year re-annotation) & --- \\
C6  & 16S identity to \emph{P.\ umbonata} = 98.38\% vs paper's 98.24\% & --- \\
C7  & Species novelty unambiguously confirmed by every method & Absolute ANI 3--7 pp above paper across three independent modern methods \\
C8--C9 & --- & AAI and dDDH not computed (heavier / web-queued); species novelty already confirmed by ANI, so not blocking \\
C10--C11 & Species novelty and phylogenetic placement replicated & --- \\
C12 & Protease claim confirmed (54 peptidases) & Chitinase (0 hits) and $\beta$-1,3-glucanase (0 hits) not supported by current PGAP annotation \\
C13--C16 & --- & All phenotypic/chemotaxonomy claims are out-of-reach from public data \\
\bottomrule
\end{longtable}

\section{Verdict + justification}

\textbf{PARTIAL} (leaning toward REPLICATED for the core taxonomic claim). This is a straightforward novel-species-announcement paper whose primary artifact is the deposited assembly. That artifact is public and every quantitative claim referring to it was reproduced within tool-drift tolerance. The species-novelty conclusion is unambiguously supported. However, several sub-claims are partially reproduced or contradicted: (a) contig count / N50 differ (NCBI cleanup), (b) the paper's own abstract-vs-body genome-length inconsistency is a genuine internal contradiction, (c) chitinase / $\beta$-1,3-glucanase claims are not supported by the current PGAP annotation, and (d) phenotypic claims cannot be tested without the live isolate. Hence PARTIAL rather than REPLICATED.

\section{Open Questions (five heavy-duty, grounded in this replication)}

\begin{enumerate}
\item Why does the paper's abstract report the genome as 2{,}476{,}083 bp while the paper's own body text and both NCBI records report 2{,}453{,}694 bp (22-kb discrepancy)? Directly observed as a paper-internal contradiction. \emph{Next step:} audit species-announcement papers 2020--2025 for abstract-vs-deposit length mismatches; if $>$5\%, propose a mandatory consistency check.
\item Why is the paper's OrthoANIu ANI (76.8\%) systematically 3--7 pp lower than fastANI (80.8\%), skani (79.4\%), and our BLASTn ANIb (83.4\%)? All three modern tools agree with each other and disagree with the paper in the same direction. \emph{Next step:} curate 30--50 published genome pairs at 75--90\% ANI where authors used OrthoANIu and raw FASTAs are public; recompute with modern tools and fit a bias regression.
\item The paper claims chitinase and $\beta$-1,3-glucanase genes (motivating a fungal-biocontrol application angle); current PGAP annotation shows zero hits for either family, despite 54 protease/peptidase hits confirming the protease claim. \emph{Next step:} re-run Prokka + dbCAN2 / dbCAN3 with explicit GH-family HMMs on \texttt{bgyt1.fna} to resolve the discrepancy.
\item The paper describes BGYT1 as ``Gram-stain-negative'' in the same document where it describes the \emph{Olsenella} genus (phylum Actinobacteria) as ``Gram stain positive'' --- Actinobacteria are canonically Gram-positive. \emph{Next step:} order KCTC 25379${}^{\mathrm{T}}$ ($\sim$US\$150), re-do Gram stain independently, and if genuine do TEM cross-sections to look for an outer membrane (which would be extraordinary).
\item The paper's closest relative \emph{O.\ umbonata} was reclassified to \emph{Parafannyhessea umbonata} (Zgheib et al.\ 2021, cited by the paper itself); should \emph{Olsenella intestinalis} therefore be reclassified as \emph{Parafannyhessea intestinalis}? \emph{Next step:} submit BGYT1 + all currently-valid \emph{Parafannyhessea}, \emph{Olsenella}, \emph{Fannyhessea}, \emph{Lancefieldella}, and \emph{Paratractidigestivibacter} type strains to TYGS for a proper GBDP tree and check whether BGYT1 clades with \emph{Parafannyhessea} or with \emph{Olsenella} sensu stricto.
\end{enumerate}

\end{document}
